%% This is file `DEMO-RWTH-SciPoster.tex' version 1.1 (2026-02-13),
%% it is part of
%% RWTH-CI -- Corporate Design for RWTH Aachen University
%% ----------------------------------------------------------------------------
%%
%% Copyright (C) 2025–2026 by Marei Peischl <marei@peitex.de>
%%
%% ============================================================================
%% This work may be distributed and/or modified under the
%% conditions of the LaTeX Project Public License, either version 1.3c
%% of this license or (at your option) any later version.
%% The latest version of this license is in
%% http://www.latex-project.org/lppl.txt
%% and version 1.3c or later is part of all distributions of LaTeX
%% version 2008/05/04 or later.
%%
%% This work has the LPPL maintenance status `maintained'.
%%
%% The current maintainer of this work is
%%   Marei Peischl <rwth-ci@peitex.de>
%%
%% The development repository can be found at
%% https://gitlab.git.nrw/rwth-it-center/rwth-latex-templates/rwth-ci
%% Please use the issue tracker for feedback!
%%
%% If you need a compiled version of this document, have a look at
%% http://mirror.ctan.org/macros/latex/contrib/rwth-ci/doc
%% or at the documentation directory of this package (if installed)
%% <path to your LaTeX distribution>/doc/latex/rwth-ci
%% ============================================================================
%%
% !TeX program = lualatex
%%

\documentclass[
english,% Main language as global option
% landscape=true,% Landscape Poster
% logo=bottom, % Move logo to bottom
% logofile=example-image-plain,% Select your logo file
% accept-missing-logos=true,% No error in case logo files are not available
]{rwth-sciposter}

%%%%%%%%%%%%%%%%%%%
% Language setup
%%%%%%%%%%%%%%%%%%%
\usepackage[english]{babel}

%%%%%%%%%%%%%%%%%%%
% Document specific setup for demonstration, generally not needed!
%%%%%%%%%%%%%%%%%%%
\newcommand*{\file}[1]{\texttt{#1}}
\newcommand*{\code}[1]{\texttt{#1}}
\newcommand*{\pkg}[1]{\textsf{#1}}
\newcommand*{\cls}[1]{\textsf{#1}}
\newcommand{\tbs}{\textbackslash}
\newcommand{\repl}[1]{<\textit{#1}>}
\newcommand*{\macro}[1]{\code{\tbs#1}}
%%%%%%%%%%%%%%%%%%%
% End of demo specific setup
%%%%%%%%%%%%%%%%%%%

\begin{document}

\title{tcolorbox-Poster using the Corporate Design of RWTH Aachen University}
\author{Marei Peischl \and the \TeX-Lion}
\footerqrcode{https://www.peitex.de}

\contactinfo{
	Telefon: +49 241 80-94123\\
	E-Mail: kim.muster@mx.rwth-aachen.de
}
\addressdata{
Institut für maximale Ausdehnung der Logosystematik\\
RWTH Aachen University\\
Musterstraße 7, 52074 Aachen, GERMANY\\
www.mx.rwth-aachen.de\\
\begin{tabular}[t]{@{}ll@{}}
Telefon:&+49 241 80-94123\\
Fax:&+49 241 80-94234
\end{tabular}
 }

\begin{tcbposter}[
poster={
rows=6,% number of rows can be changed. 6 is the default
%showframe,% Might be helpful for placements
}]

\begin{posterboxenv}[title=Intro]{name=intro,column=1,row=1,span=2}
	The document class \cls{rwth-sciposter} is based on the \pkg{tcolorbox} package by Thomas F. Sturm.
	This document is an example for an academic poster using the Corporate Design of RWTH Aachen University and may be used as a template file.

\end{posterboxenv}

% Box with an image without title but using a caption
\begin{posterboxenv}{name=caption,below=intro,span=3,rowspan=2}
	\includegraphics[width=\linewidth]{example-image-plain}
	\captionof{figure}{
		An example image in a box with a caption
	}
\end{posterboxenv}


\begin{posterboxenv}[title=Placement of boxes]{name=positioning,column=4,span=2}
	It is an additional step, but the boxes are positioned manually in the \pkg{poster} library of the \pkg{tcolorbox} package.
	This not only allows a nicer alignment of the boxes, but also simplifies the positioning of cross-references.
	See also \pkg{tcolorbox} documentation.
\end{posterboxenv}

\begin{posterboxenv}[title=QR code]{name=qrcode,column=4,span=3,row=5}
	The QR code is set using \macro{footerqrcode}.

	This setup also provides a mechanism to use a URL set via \pkg{biblatex}.
	The variant \macro{footerqrcode*} will use the argument as a citation key and create a QR code for the corresponding URL.
\end{posterboxenv}



\begin{posterboxenv}[title=Footer]{name=footer,below=qrcode,column=4,span=3}
	The data used in the footer can be set using the custom commands \macro{contactinfo} and \macro{addressdata}.

	They both use one argument which will be placed within the corresponding box in the footer.
\end{posterboxenv}

\begin{posterboxenv}[title=Relative positioning]{name=relative,between=positioning and qrcode,column=4,span=3}
	This box is placed between \code{positioning} box and \code{qrcode} box.
	The box is placed on top of the free area.

	Of course this code has to be placed after the boxes which are used for reference.

	It's possible to combine multiple parameters for exact placement. Please have a look at the tcolorbox documentation for details.
\end{posterboxenv}



\end{tcbposter}

\end{document}
%% End of file `DEMO-RWTH-SciPoster.tex'.
